\begin{textsample}{POS Dim 4 – 2010\_c6 – Score -1.00 – 2010\_c6/enade\_2010\_med\_c0c86d6464.txt}  \label{ex:f4_pos_313}
CAPÍTULO 3 ANÁLISE TÉCNICA DA PROVA Este capítulo tem por objetivo apresentar o desempenho dos estudantes de Medicina no ENADE/2010.
Para isso, foram calculadas as estatísticas básicas da prova em seu todo, bem como as estatísticas dos componentes relacionadas à Formação Geral e ao de Conhecimento Específico da Área.
Nas tabelas, são apresentados o tamanho da população inscrita e de presentes, e as seguintes estatísticas das notas7 : média do desempenho na prova, erro padrão da média, desvio padrão, nota mínima, mediana, nota máxima, coeficiente de assimetria e de variação.
As estatísticas apresentadas neste capítulo contemplam o total de estudantes e, separadamente, os Ingressantes e os Concluintes.
Foram calculadas tendo -se em vista as seguintes agregações : ( a ) as Grandes Regiões e o país como um todo ; ( b ) a Categoria Administrativa ; e ( c ) a Organização Acadêmica.
Em relação aos gráficos de distribuição de notas, o intervalo considerado foi de 10 unidades, aberto à esquerda e fechado à direita, com exceção do primeiro intervalo, [ 0 ; 10 ], fechado em ambos os extremos.
Para os gráficos de distribuição das notas das questões discursivas, foram consideradas mais duas categorias : questão em branco e nota zero. 3.1 ESTATÍSTICAS BÁSICAS DA PROVA 3.1.1 Estatísticas Básicas Gerais A Tabela 3.1 apresenta as estatísticas básicas da prova por grupo de estudantes.
A população total de estudantes regularmente inscritos8 para a prova foi de 30.004.
Destes, 28.938 estiveram presentes, sendo de 3,5 \% o índice de não comparecimento.
A abstenção foi maior entre Ingressantes ( 5,0 \% ) do que entre Concluintes ( 1,2 \% ).
A média de toda a prova foi 47,4, sendo que os Ingressantes obtiveram média mais baixa, 40,0, que os Concluintes, cuja média foi 58,8.
O desvio padrão foi 15,0, sendo o do grupo de Ingressantes menor ( 10,8 ) do que o do grupo dos Concluintes ( 13,4 ) ; os coeficientes de variação ( 27,0 \% no grupo dos Ingressantes e 22,8 \% no dos Concluintes ) indicam menor dispersão relativa nas notas de Concluintes.
A nota máxima foi 90,9, obtida por Concluintes, ao passo que a maior nota obtida por Ingressantes foi 84,8. 7 Algumas dessas estatísticas estão definidas no Capítulo 1. 8 As estatísticas das notas apresentadas consideram apenas os estudantes regularmente inscritos e com resultados validados por ausência de erro administrativo.
Uma comparação entre o desempenho de Ingressantes e Concluintes também pode ser realizada a partir dos dados apresentados no Gráfico 3.1, com a distribuição de notas finais.
Destaca -se que as notas mais elevadas, considerando -se o intervalo posterior à nota 50,0, foram atingidas fundamentalmente pelos Concluintes, cuja concentração das notas está à direita dos Ingressantes.
Já nas faixas de notas mais baixas, inverteu -se o quadro identificado nas faixas de notas mais altas.
Essa característica é corroborada pelo valor da nota mediana 40,8 no grupo dos Ingressantes, abaixo da encontrada no grupo dos Concluintes, 60,0, transparecendo uma superioridade de desempenho dos Concluintes em relação aos Ingressantes na prova em tela ( Tabela 3.1 ).
Os coeficientes de assimetria indicam uma assimetria negativa ( -0,8 ) para os Concluintes, bem como para Ingressantes ( -0,9 ).
O Gráfico 3.2 apresenta informações referentes à nota final de Ingressantes e Concluintes, comparando -se os resultados de acordo com as Grandes Regiões do país, a Categoria Administrativa e a Organização Acadêmica.
No que diz respeito às notas médias dos alunos em cada Grande Região, observa -se que, em relação às notas aferidas em âmbito nacional ( média de 40,0 para Ingressantes e de 58,8 para Concluintes ), existe diferença estatisticamente significativa ao nível de 95 \% entre a menor média, obtida na região Sudeste ( 38,7 ) e a maior, na região Nordeste ( 42,8 ), no grupo de Ingressantes.
Entre os Concluintes, também existe diferença estatisticamente significativa ao nível de 95 \% das médias.
A maior média de notas foi obtida pela região Sul ( 61,7 ), ao passo que a menor foi identificada na região Norte ( 54,9 ).
As demais médias de notas foram 57,6 para a região Sudeste, 60,4 para a região Nordeste e 60,7 para a região Centro-Oeste.
Destaca -se a região Sul com médias superiores às nacionais, tanto para o grupo de Ingressantes quanto para o de Concluintes.
Levando -se em conta os grupos dos estudantes em Categorias Administrativas, observa -se que a pontuação média mais elevada entre os Concluintes foi encontrada em instituições públicas ( 62,0 ), com desempenho acima da média nacional.
Já os Concluintes provenientes de instituições privadas obtiveram média 56,1 \%, abaixo da média nacional.
Quanto aos Ingressantes, as notas mais elevadas ocorreram também nas instituições públicas ( 40,7 ).
Os Ingressantes vinculados a instituições privadas obtiveram nota média 39,5, abaixo da média nacional.
Tendo como foco as notas médias dos estudantes agrupados por Organização Acadêmica, distribuídas entre Universidades, Centros Universitários e Faculdades, para os Concluintes vinculados a Universidades, a média 59,6 ficou acima da média nacional.
Os Ingressantes de Universidades tiveram nota média, 40,1, um pouco acima da nacional, mas não estatisticamente diferente a 95 \%.
Considerando -se as demais Organizações Acadêmicas, observa -se que nos Centros Universitários e Faculdades, o desempenho de Concluintes ficou abaixo da média nacional.
Concluintes de Centros Universitários e de Faculdades tiveram notas médias 53,1 e 57,8, respectivamente.
Ingressantes das mesmas Organizações Acadêmicas obtiveram notas 40,0 e 39,8, respectivamente, não estatisticamente inferiores ( a 95 \% ) da média nacional. 3.1.2 Estatísticas Básicas no Componente de Formação Geral A Tabela 3.2 apresenta as estatísticas básicas em relação ao componente da prova que avalia a Formação Geral dos estudantes.
Os Concluintes obtiveram desempenho médio 58,3, um pouco superior ao dos Ingressantes, 58,0.
Quanto à variabilidade das notas nos dois grupos, o desvio padrão das notas obtidas foi 17,1.
Os dois grupos alcançaram notas máximas iguais : 100,0.
O Gráfico 3.3 propicia a comparação entre Concluintes e Ingressantes relativamente ao desempenho no componente de Formação Geral.
Nesse componente, os Concluintes e Ingressantes apresentaram desempenho semelhante.
Dos Ingressantes e Concluintes, um pouco mais de 26,0 \% obtiveram notas na faixa entre 60 e 70 pontos, sendo essa a faixa modal para ambos os grupos Nas faixas de notas seguintes, mais elevadas, verifica -se queda gradual de ocorrência de estudantes.
Entre os Concluintes, houve percentual maior de estudantes ( em comparação com os Ingressantes ) em todos os intervalos anteriores à faixa de 60 a 70 pontos, com exceção do primeiro intervalo : [ 0 ; 10 ], onde a proporção de Ingressantes é maior.
Apenas 2,5 \% dos Ingressantes e 1,2 \% dos Concluintes tiveram desempenho nesta faixa de notas.
O coeficiente de assimetria foi negativo : -0,7 para Concluintes e -0,9 para Ingressantes.
No Gráfico 3.4 são apresentadas as informações referentes ao desempenho dos Ingressantes e dos Concluintes no componente de Formação Geral, em diferentes agregações como Grande Região do país, Categoria Administrativa e Organização Acadêmica.
Considerando -se as notas médias dos estudantes segundo cada Grande Região, observa -se que no grupo de Ingressantes existe diferença estatisticamente significativa ao nível de 95 \% entre as médias na região Nordeste ( 60,8, a mais alta ) e Norte ( 56,6, a mais baixa ).
No grupo de Concluintes, as notas mais elevadas ocorreram nas regiões Nordeste ( 61,4 ), seguida pela Centro-Oeste ( 61,2 ).
A diferença entre notas mais altas e mais baixas para Concluintes foi também estatisticamente significativa ao nível de 95 \% : a menor nota foi 56,4, correspondente aos da região Sudeste.
Essa região ficou abaixo da média nacional no componente de Formação Geral, com nota 56,9 para Ingressantes e 56,4 para Concluintes.
Tendo como foco as Categorias Administrativas, as instituições públicas obtiveram médias mais elevadas do que as privadas : para Concluintes, notas médias de 61,1 em cursos de IES públicas e de 56,1 nos das privadas.
Para os Ingressantes, as notas foram de, respectivamente, 59,9 e 56,8.
Em comparação à média nacional, as notas de Ingressantes e de Concluintes vinculados às instituições públicas foram superiores, ao passo que, nas instituições privadas, foram inferiores.
Considerando -se o tipo de Organização Acadêmica, nota -se também uma diferença estatisticamente significativa a 95 \% entre a maior e a menor média.
Nas Universidades, houve desempenho melhor do que em Centros Universitários e Faculdades.
Neste primeiro tipo de Organização Acadêmica, Ingressantes tiveram nota média 58,5, enquanto Concluintes, 58,9.
Nos Centros Universitários e Faculdades, os valores foram mais baixos : os Ingressantes tiveram média 56,9 nos Centros Universitários e 57,2 nas Faculdades ; entre Concluintes, a média foi 54,2 para Centros Universitários e 57,7 para as Faculdades.
Observa -se que a nota de Ingressantes foi superior a de Concluintes nos Centros Universitários. 3.1.3 Estatísticas Básicas do Componente de Conhecimento Específico A Tabela 3.3 apresenta as estatísticas básicas referentes à parte do componente de Conhecimento Específico da Área de Medicina.
A média do desempenho foi 43,8.
Os Concluintes apresentaram desempenho médio ( 58,9 ) superior ao dos Ingressantes ( 34,0 ) ; as medidas de dispersão – desvio padrão e amplitude geral das notas – são menores no grupo dos Ingressantes do que no grupo dos Concluintes.
Os coeficientes de variação, 32,6 \% no grupo dos Ingressantes e 25,5 \% no grupo dos Concluintes, indicam dispersão relativa menor nas notas desses últimos.
As notas máximas obtidas por Concluintes e Ingressantes foram iguais a 96,0 e 92,0, respectivamente, notas inferiores àquelas encontradas no componente de Formação Geral ( 100,0 e 100,0, respectivamente ).
As medianas das distribuições das notas dos Ingressantes e dos Concluintes foram 34,4 e 59,6, respectivamente, enfatizando a superioridade dos Concluintes neste componente da prova.
Assim como os Gráficos 3.1 e 3.3, o Gráfico 3.5, apresentado a seguir, proporciona uma comparação do desempenho entre Concluintes e Ingressantes, em relação ao componente de Conhecimento Específico.
Dentre as três comparações feitas, observa -se que, embora a distribuição das notas seja semelhante à apresentada na prova como um todo, neste componente, a diferença de desempenho entre Concluintes e Ingressantes é mais acentuada, destacando -se o desempenho mais elevado dos primeiros nas faixas de notas mais altas.
Claramente a distribuição das notas dos Ingressantes, mais concentradas à esquerda em relação aos Concluintes, evidencia o percentual mais alto de notas nas faixas mais baixas.
A maior concentração de Concluintes está nos intervalos compreendidos entre ( 50 a 60 ] e ( 60 a 70 ] pontos ( sendo o primeiro intervalo, o modal ), ao passo que a de Ingressantes está entre ( 20 a 30 ] e ( 30 a 40 ] pontos ( sendo o segundo intervalo, o modal ).
Assim como os Gráficos 3.2 e 3.4, o Gráfico 3.6 apresenta uma comparação dos resultados em relação às Grandes Regiões do país, à Categoria Administrativa e à Organização Acadêmica, agora levando em conta o desempenho dos Ingressantes e dos Concluintes no componente de Conhecimento Específico da prova.
Ao se observar as notas médias dos estudantes, segundo cada Grande Região do país, verifica -se que, entre os Concluintes, a mais elevada foi encontrada na região Sul ( 62,2 ), maior que a média nacional ( 58,9 ).
Em relação aos Ingressantes, a região Nordeste apresenta a média mais elevada ( 36,8 ), seguida da região Sul, com média igual a 34,5.
A média mais baixa entre os Concluintes ( 54,0 ) foi verificada na região Norte, e entre os Ingressantes ( 32,6 ), na Sudeste.
Quanto à Categoria Administrativa, observa -se comportamento semelhante àquele da parte relativa à Formação Geral e à prova como um todo, ou seja, as médias mais elevadas tanto entre os Concluintes quanto entre os Ingressantes foram encontradas em instituições de ensino públicas.
Nesse tipo de IES, a nota média de Concluintes foi 62,3, e a de Ingressantes, 34,2.
Tais médias estão situadas acima da média nacional.
Nas instituições privadas, as médias foram, mais uma vez, inferiores em relação às públicas.
Concluintes obtiveram média 56,1 e Ingressantes, 33,8.
Em relação à Organização Acadêmica, o desempenho dos Concluintes foi superior ao dos Ingressantes nas três categorias.
Nas Universidades, os Concluintes obtiveram a média mais elevada, igual a 59,9.
Um pouco abaixo estiveram as Faculdades ( 57,8 ) e os Os Ingressantes dos três tipos de Organização Acadêmica tiveram notas semelhantes, sem nenhuma diferença estatisticamente significativa : os de Universidades atingiram 33,9, os de Centros Universitários, 34,3, e os de Faculdades, 34,0. 3.2 ANÁLISE DAS QUESTÕES OBJETIVAS 3.2.1 Componente de Formação Geral A Tabela 3.4 apresenta as estatísticas básicas relativas às oito questões objetivas do componente da prova que abrange a Formação Geral dos estudantes.
Como pode ser observado, os estudantes Concluintes obtiveram desempenho um pouco acima, com nota média 62,9, em relação aos alunos Ingressantes, cuja nota média foi 61,2.
As medidas de dispersão nos dois grupos de alunos foram similares, como indicam os desvios padrões, assim como a amplitude total das notas, sendo que a nota mínima foi zero e a máxima, 100,0, para os dois grupos de alunos avaliados.
Os coeficientes de variação, 30,4 \% no grupo dos Ingressantes e 27,0 \% no grupo dos Concluintes, indicam menor dispersão relativa nas notas desses últimos.
A mediana foi igual a 62,5 nos dois grupos.
O coeficiente de assimetria do total de estudantes revelou assimetria negativa ( -0,8 ).
A Tabela 3.5 apresenta as classificações das questões objetivas do componente de Formação Geral, segundo o índice de facilidade.
As questões objetivas da prova foram assim avaliadas : das oito questões, duas tiveram o índice de facilidade classificado como muito fácil.
Duas questões foram tidas como fáceis, por terem índice de acertos situado na faixa entre 0,61 e 0,85 ( de 61,0 \% a 85,0 \% de acertos ).
Uma questão foi considerada de dificuldade média, situando -se no intervalo entre 0,41 e 0,60 do índice de facilidade, ou seja, houve entre 41,0 \% e 60,0 \% de acertos.
Três questões foram classificadas como difíceis, situando -se no intervalo entre 0,16 e 0,40.
Por fim, nenhuma questão foi classificada como muito difícil.
A Tabela 3.6 descreve os resultados das análises das questões objetivas relativas à Formação Geral, segundo o poder de discriminação, utilizando -se para tal, do índice de discriminação ponto bisserial.
Sete das oito questões apresentaram índices iguais ou acima de 0,40 e, assim, foram classificadas como muito boas nesta característica para esse grupo de alunos ; uma questão teve bom índice de discriminação, entre 0,30 e 0,39.
Nenhuma das questões do Componente de Formação Geral foi classificada como de média e de fraca discriminação.
Na sequência, na Tabela 3.7, estão especificadas, por questão, os resultados da análise relativa aos índices de facilidade e de discriminação de cada questão no componente relativo à Formação Geral da prova de Medicina.
O índice de facilidade variou de 0,17 a 0,94, e o de discriminação, de 0,33 a 0,51.
Das questões com índices de discriminação muito bom, de números 1, 2, 3, 4, 5, 6 e 8, somente duas delas ( questões 3 e 8 ) figuraram como muito fácil nesse conjunto.
As questões 5 e 8 foram as que apresentaram maior poder discriminatório, ambas com índice 0,51.
A questão de número 7 apresentou índice de discriminação 0,33, bom, mas foi classificada no índice de facilidade como difícil, com 0,17, ou seja, um quantitativo pequeno de estudantes conseguiu resolvê -la, dentro do universo de participantes.
A Tabela 3.6 descreve os resultados das análises das questões objetivas relativas à Formação Geral, segundo o poder de discriminação, utilizando -se para tal, do índice de discriminação ponto bisserial.
Sete das oito questões apresentaram índices iguais ou acima de 0,40 e, assim, foram classificadas como muito boas nesta característica para esse grupo de alunos ; uma questão teve bom índice de discriminação, entre 0,30 e 0,39.
Nenhuma das questões do Componente de Formação Geral foi classificada como de média e de fraca discriminação.
Na sequência, na Tabela 3.7, estão especificadas, por questão, os resultados da análise relativa aos índices de facilidade e de discriminação de cada questão no componente relativo à Formação Geral da prova de Medicina.
O índice de facilidade variou de 0,17 a 0,94, e o de discriminação, de 0,33 a 0,51.
Das questões com índices de discriminação muito bom, de números 1, 2, 3, 4, 5, 6 e 8, somente duas delas ( questões 3 e 8 ) figuraram como muito fácil nesse conjunto.
As questões 5 e 8 foram as que apresentaram maior poder discriminatório, ambas com índice 0,51.
A questão de número 7 apresentou índice de discriminação 0,33, bom, mas foi classificada no índice de facilidade como difícil, com 0,17, ou seja, um quantitativo pequeno de estudantes conseguiu resolvê -la, dentro do universo de participantes.
O Gráfico 3.7 analisa o comportamento da questão de número 8 de Formação Geral.
Trata -se da questão mais fácil e com o maior índice de discriminação desse componente da prova.
No gráfico a seguir, cada uma das cinco curvas representa o percentual de respostas em determinada alternativa da questão, em função da nota dos estudantes nesta parte da prova ( Formação Geral/Múltipla Escolha ) antes de possíveis eliminações pelo critério do ponto bisserial.
A curva em rosa corresponde à alternativa B, a correta para este item.
Assim, observa -se que entre os estudantes com notas mais baixas, desempenho mais fraco nesse componente do exame, a situação mais frequente foi a escolha de uma das alternativas incorretas : a alternativa C ( em verde ) ou E ( em roxo ).
Na medida em que a nota aumenta, indicando melhor desempenho nessa parte da prova, aumenta concomitantemente a proporção de estudantes que selecionaram a alternativa correta B, atingindo 100 \% para as notas acima de 70 pontos.
Essa análise permite verificar como a questão discriminou entre os grupos de desempenho, justificando o alto índice obtido na questão.
Os gráficos relativos às demais questões de Formação Geral constam do Anexo I. 3.2.2 Componente de Conhecimento Específico A Tabela 3.8 apresenta as estatísticas básicas em relação às questões objetivas do componente de Conhecimento Específico da prova.
Nesse conjunto de questões, os estudantes Concluintes ( 61,8 ) obtiveram nota média superior, aos Ingressantes ( 34,1 ).
Quanto à variabilidade das notas, o desvio padrão no grupo dos Ingressantes foi 11,8, e no grupo dos Concluintes, 16,6.
Os coeficientes de variação foram 34,6 \% e 26,9 \%, respectivamente, indicando menor dispersão relativa nas notas dos Concluintes.
Com relação à mediana, no grupo dos Ingressantes, esta foi igual a 36,4, estando acima da média desse grupo, enquanto no grupo dos Concluintes, a mediana foi igual a 63,6, também superando a média do grupo.
O coeficiente de assimetria foi ( -0,2 ) para Ingressantes e ( -0,5 ) para Concluintes.
A Tabela 3.9 apresenta as classificações das questões objetivas da parte da prova relativa ao componente de Conhecimento Específico, segundo o índice de facilidade.
A partir dos índices obtidos, pode -se concluir que boa parte das questões objetivas da prova foram tidas como difíceis : das 27 questões, treze foram classificadas como difíceis e uma como muito difícil.
Não houve questão classificada como muito fácil, ao passo que cinco foram tidas como fáceis, na faixa de 0,61 a 0,85 do índice de facilidade, enquanto oito foram consideradas médias, entre 0,41 e 0,60.
Na Tabela 3.10 são apresentados os índices de discriminação das questões objetivas do componente de Conhecimento Específico da prova.
Observa -se, segundo os índices estabelecidos, que a capacidade de discriminação das questões objetivas desse componente foi classificada, para cinco das 27 questões, como boa, enquanto quatorze delas tiveram índice de discriminação muito bom.
A maioria das questões – 19 em 27 – proporcionou, portanto, índices de discriminação bons ou muito bons.
Três delas foram classificadas como médias e cinco como fracas, sendo oito, por conseguinte, a quantidade de questões nos dois patamares mais baixos de discriminação.
Constata -se, assim, que a A Tabela 3.9 apresenta as classificações das questões objetivas da parte da prova relativa ao componente de Conhecimento Específico, segundo o índice de facilidade.
A partir dos índices obtidos, pode -se concluir que boa parte das questões objetivas da prova foram tidas como difíceis : das 27 questões, treze foram classificadas como difíceis e uma como muito difícil.
Não houve questão classificada como muito fácil, ao passo que cinco foram tidas como fáceis, na faixa de 0,61 a 0,85 do índice de facilidade, enquanto oito foram consideradas médias, entre 0,41 e 0,60.
Na Tabela 3.10 são apresentados os índices de discriminação das questões objetivas do componente de Conhecimento Específico da prova.
Observa -se, segundo os índices estabelecidos, que a capacidade de discriminação das questões objetivas desse componente foi classificada, para cinco das 27 questões, como boa, enquanto quatorze delas tiveram índice de discriminação muito bom.
A maioria das questões – 19 em 27 – proporcionou, portanto, índices de discriminação bons ou muito bons.
Três delas foram classificadas como médias e cinco como fracas, sendo oito, por conseguinte, a quantidade de questões nos dois patamares mais baixos de discriminação.
Constata -se, assim, que a Na Tabela 3.11 estão especificados, por questão, os resultados da análise relativa aos índices de facilidade e de discriminação de cada uma das questões do componente de Conhecimento Específico da prova de Medicina.
Dentre as questões que alcançaram os maiores índices de discriminação, as de números 13, 14, 16, 17, 19, 20, 25, 27, 28, 29, 32, 33, 34 e 35 foram classificadas com o índice muito bom, situando -se no intervalo acima de 0,40.
Seus índices de discriminação situaram -se de 0,41 a 0,52.
A questão de número 37 foi a mais difícil dentre as 27 questões específicas válidas, com baixo índice de facilidade, 0,14, o que, em termos percentuais, corresponde a 14,0 \% de acertos.
Essa questão apresentou poder discriminatório igualmente baixo, 0,14, o que indica ter sido esta uma das mais difíceis para os estudantes, mas foi eliminada do cômputo da nota final pelo critério.
Destacam -se, também, as questões de números 22 e 34, para as quais os índices de facilidade foi 0,18 em ambas, o que, em termos percentuais, corresponde a 18,0 \% de estudantes que responderam acertadamente.
Tais questões ficaram portanto,entre as mais difíceis da prova.
Ainda a notar as questões que obtiveram índice de discriminação fraco, -0,07 para a 12 e 0,10 para a questão 36. prova – no que se refere ao componente de Conhecimento Específico – possuía boa capacidade de discriminar entre aqueles que dominam ou não o conteúdo, ainda que cinco das questões tenham sido anuladas pelo critério e não tenham sido utilizadas para o cálculo das notas finais : questões 12, 15, 30, 36 e 37.
Na Tabela 3.11 estão especificados, por questão, os resultados da análise relativa aos índices de facilidade e de discriminação de cada uma das questões do componente de Conhecimento Específico da prova de Medicina.
Dentre as questões que alcançaram os maiores índices de discriminação, as de números 13, 14, 16, 17, 19, 20, 25, 27, 28, 29, 32, 33, 34 e 35 foram classificadas com o índice muito bom, situando -se no intervalo acima de 0,40.
Seus índices de discriminação situaram -se de 0,41 a 0,52.
A questão de número 37 foi a mais difícil dentre as 27 questões específicas válidas, com baixo índice de facilidade, 0,14, o que, em termos percentuais, corresponde a 14,0 \% de acertos.
Essa questão apresentou poder discriminatório igualmente baixo, 0,14, o que indica ter sido esta uma das mais difíceis para os estudantes, mas foi eliminada do cômputo da nota final pelo critério.
Destacam -se, também, as questões de números 22 e 34, para as quais os índices de facilidade foi 0,18 em ambas, o que, em termos percentuais, corresponde a 18,0 \% de estudantes que responderam acertadamente.
Tais questões ficaram portanto, entre as mais difíceis da prova.
Ainda a notar as questões que obtiveram índice de discriminação fraco, -0,07 para a 12 e 0,10 para a questão 36.
O Gráfico 3.8 analisa o comportamento da questão 29 do componente de Conhecimento Específico.
Essa questão teve índice de facilidade 0,33, ou seja, 33,0 \% dos estudantes assinalaram acertadamente a opção E, correspondente ao gabarito.
Seu índice de discriminação foi igual a 0,44, classificado como muito bom, embora outras questões que apresentaram percentuais mais elevados de respostas corretas tenham tido índices de discriminação mais elevados, como o caso da questão 19, com 0,47 de índice de facilidade e 0,52 de índice de discriminação.
Neste gráfico, cada uma das cinco curvas representa o percentual de respostas em determinada alternativa da questão 29, em função da nota dos estudantes nesta parte da prova antes de possíveis eliminações de questões pelo critério do ponto bisserial.
A alternativa correta E, representada no gráfico pela curva em roxo, foi escolhida em maiores proporções pelos alunos com desempenho melhor, enquanto as alternativas incorretas, também denominadas distratores, foram selecionadas principalmente por aqueles com notas mais baixas, ainda que aqueles com nota zero, na sua quase totalidade tenham deixado em branco ou marcaram mais de uma alternativa, comportamento considerado inválido.
A proporção de alunos que selecionaram a resposta correta E aumenta gradativamente com a nota dos alunos chegando a atingir 100 \% para notas acima de 90, enquanto a proporção dos que escolhem alternativas incorretas, decai a partir da nota 20 como função da nota nesta parte da prova.
Os gráficos relativos às demais questões do componente de Conhecimento Específico constam do Anexo I. 3.3 ANÁLISE DAS QUESTÕES DISCURSIVAS 3.3.1 Componente de Formação Geral A análise dos resultados de desempenho dos estudantes de Medicina, nas duas questões discursivas relativas à Formação Geral, encontra -se na Tabela 3.12 e no Gráfico 3.9.
Observa -se, a partir destes, que as notas médias foram mais baixas nesse conjunto que no das objetivas.
A média geral dos Ingressantes nas questões objetivas de Formação Geral foi 61,2 ( Tabela 3.4 ), ao passo que nas questões discursivas, a média foi 62,9.
O mesmo ocorreu entre os Concluintes, que tiveram, em Formação Geral, média 58,3 nas questões objetivas e 51,4 nas questões discursivas.
Cabe destacar que, nas questões discursivas, o desempenho dos Ingressantes foi ligeiramente superior ao desempenho dos Concluintes.
Os dois agrupamentos atingiram a mesma nota máxima, 100,0, o que acompanha a constatação de que não foi muito evidente uma superioridade dos Concluintes, em relação aos Ingressantes, nesta parte do exame.
Os coeficientes de assimetria indicam assimetria negativa para Ingressantes e Concluintes.
Verifica -se a semelhança das distribuições das notas desses grupos, no Gráfico 3.9, apresentado a seguir.
Na sequência, os resultados verificados para cada uma das questões discursivas de Formação Geral serão apresentados, estabelecendo -se relações com os conteúdos englobados em cada uma delas.
Os comentários da banca corretora a respeito do observado nas respostas dos estudantes, suas impressões e conclusões serão expostos, junto à análise de cada questão.
Cumpre esclarecer que, tendo em vista que as questões discursivas de Formação Geral são padronizadas, ou seja, constam de todas as provas, os comentários da banca são os mesmos para todas as carreiras acadêmicas, sendo direcionados a todos os estudantes que participaram do ENADE/2010.
A seguir, serão analisados os desempenhos de Ingressantes e Concluintes da Área de Medicina nas duas questões discursivas de Formação Geral do ENADE/2010. 3.3.1.1 Análise da Questão Discursiva 9 do Componente de Formação Geral Os dados obtidos a partir das respostas à questão 9 encontram -se na Tabela 3.13 e no Gráfico 3.10.
Nessa questão – de melhor desempenho dentre as duas de Formação Geral – os Ingressantes tiveram média 57,0, e os Concluintes, 54,0, resultando com diferença de 3,0 pontos entre os dois grupos.
Houve semelhança na distribuição de notas dos dois grupos, observando -se percentuais aproximados de Concluintes e Ingressantes que deixaram a questão em branco, 12,9 \% dos Ingressantes e 13,7 \% dos Concluintes, enquanto obtiveram nota zero na questão mesmo sem tê -la deixado em branco, 1,6 \% e 1,7 \%, respectivamente.
Entre Ingressantes, a nota mediana foi 65,0 pontos, enquanto Concluintes obtiveram 60,0.
Ambos os grupos tiveram como modal o intervalo de 60 a 70 pontos, com respectivamente 18,5 e 18,1 \% dos estudantes Ingressantes e Concluintes nessa faixa.
A nota máxima, 100,0 pontos, foi alcançada tanto pelos Concluintes quanto pelos Ingressantes. 3.3.1.2 Comentários sobre as respostas à Questão Discursiva 9 No que tange mais precisamente à atualidade do tema proposto e ao desenvolvimento do conteúdo, deve ser ressaltado que se mostrou sintonizado com a realidade política brasileira, bem como com a necessária reflexão sobre o papel na sociedade do profissional formado pelas Instituições de Ensino Superior do país.
Dissertar sobre democracia e ética foi algo apropriado para a época em que a prova foi aplicada, pois coincidiu com o período das eleições, em que tais temas são mais discutidos na vida cotidiana.
Não se pode, entretanto, deixar de se levar em conta a tendência de os estudantes terem limitado a associação do tema “ ética ” ao assunto “ política ”, considerando -se que a ética abarca não apenas a política.
Tal associação atingiu, em alguma medida, o alcance mais preciso das respostas.
As respostas apresentadas refletiram motivação para escrever sobre o tema da questão.
Alguns estudantes chegaram a desenvolver textos eloquentes e engajados.
Em alguns casos, o conceito de democracia, por exemplo, foi bem desenvolvido ao longo do texto, simultaneamente às exemplificações pedidas.
Além disso, muitos estudantes associaram o exercício da ética ao seu futuro ofício, reconhecendo o seu papel como cidadão no desempenho de sua profissão.
Essa associação insinuou, também, o esforço de reconhecimento da noção de ética prática, ainda que de maneira codificada.
Um dos pontos positivos revelou -se no fato de, na maioria das respostas analisadas, haver uma preocupação em definir a sociedade democrática.
Isso facilitou a correção e elevou a média de notas.
Outro aspecto positivo foi a presença pertinente de exposições coerentes de comportamento ético de um futuro profissional.
A articulação dos itens “ a ” e “ b ” da questão produziu, portanto, todo tipo de reflexão acerca de uma postura ética na sociedade, inclusive, corporativa.
A comprovação de que o tema ética tem sido abordado nos cursos superiores foi outro destaque positivo detectado durante as correções.
Assim, além do caráter crítico de que muitas das respostas corretas se revestiram, os relatos constantes sobre desdobramentos da ética comprovaram a abordagem sobre o tema, no mínimo, a ética profissional, nas carreiras universitárias que participaram do ENADE.
A despeito de alguns estudantes se posicionarem contra o exame, seja anulando a questão ou escrevendo algum comentário desabonador sobre o processo, foi observado que significativa parte dos participantes tem consciência político-acadêmica e percebe por que está realizando aquele tipo de prova.
Desse modo, a maioria dos estudantes tinha algo a “ dizer “ sobre o tema da questão, isto é, valia -se da oportunidade para se posicionar a respeito do assunto trazido à baila.
Ainda do ponto de vista do conteúdo desenvolvido nas respostas, observou -se preocupação solidária e consciência cidadã dos estudantes.
Mais um destaque positivo crucial se refere à habilidade textual.
Os estudantes, por exemplo, expressaram atitudes éticas e não éticas e as contrapunham ao conceito de democracia.
Há casos em que os alunos definiam a corrupção política como um exemplo de atitude não ética, enfatizando a real missão desse cargo em uma sociedade democrática : “ atender aos anseios da população e zelar pelo bem comum ”, escreveram vários estudantes.
No tocante à produção de texto dissertativo, à objetividade e ao caráter sintético, muitos estudantes corresponderam ao solicitado, elaborando pequenos textos dissertativos.
Nessas respostas, percebeu -se demonstração de domínio do assunto focado e utilização dos textos de referência como base para a elaboração.
O enunciado da questão 9 determinou que fosse elaborado um texto dissertativo que contivesse a )..., b )... e c )..., fornecendo um encaminhamento para as respostas que conduziu tanto a um resultado positivo quanto a um negativo, a saber : foi produtiva a resposta dos estudantes no sentido de que conseguiram, na maior parte das vezes, atender plenamente aos três quesitos apresentados na questão ; por outro lado, mais especificamente devido ao uso das três letras indicativas ( a, b e c ), houve tendência à elaboração de respostas divididas em itens, em vez de um texto eminentemente dissertativo.
Muitos estudantes apresentaram respostas nesse formato fragmentado para satisfazer a formalidade do roteiro sugerido no comando da questão.
A ideia-chave que deveria permear toda a resposta – democracia – se perdia em fragmentos desconexos entre si.
Se, por um lado, o roteiro do comando da questão facilitou o raciocínio de alguns dos estudantes, por outro, induziu grande parte deles a dois caminhos : i ) o das respostas curtas e objetivas, do tipo “ um profissional de saúde não ético é aquele que conta para as outras pessoas o que foi feito na consulta com o seu paciente ” e ii ) o outro caminho, que é do contraponto entre o conceito de sociedade democrática e a realidade brasileira.
Esse segundo caminho levou vários universitários a produzirem textos bem escritos, distantes, entretanto, do tema exigido.
No que se reporta às especificidades do desenvolvimento do conteúdo, a dificuldade ou falta de competência comunicativa revelada em grande parte das respostas denuncia a pouca exposição dos alunos à leitura de textos escritos de mais excelência.
Tal dificuldade mostrou -se mais evidente na impropriedade no uso de conectores e na incapacidade de elaboração de texto dissertativo, no tratamento de questões mais abstratas.
Muitos estudantes demonstraram não saber em que consistia a ética e confundiram exemplificação com definição : ao tentarem dar exemplos, via de regra, elaboravam uma definição incipiente do que se pedia.
Portanto, a dificuldade dos estudantes em argumentar e articular as ideias com coesão e coerência ainda são os grandes “ vilões ” desse processo avaliativo.
O excesso de fragmentos e textos prévios acerca do tema permitiu que a elaboração das respostas se resumisse a uma colagem ou paráfrase mal feita dos materiais disponíveis.
Em vez de subsidiar as respostas, esse material acabou por influenciar negativamente o desempenho dos estudantes.
Desse modo, verificou -se um nível de letramento ainda significativamente baixo, visto que a maioria apresentou uma visão bastante reducionista e simplificada sobre o tema em pauta.
De um modo geral, ao conceituarem ética e sociedade democrática, os estudantes demonstraram desconhecer o limite entre os dois conceitos, sem aludirem à interseção de ambos.
Foi expressiva a quantidade de redações nas quais questionavam a democracia pelo simples fato de estarem fazendo, a fortiori, aquela prova especificamente.
Em muitas das respostas, os estudantes se diziam incomodados pelo fato de estarem obrigados a realizar o exame, demonstrando sua insatisfação.
Por diversas vezes, os estudantes não foram capazes de sintetizar e de selecionar seu discurso o suficiente para atender ao que era pedido, o que ocasionou a falta de, pelo menos, um dos aspectos da resposta.
Foi constatada grande parte de respostas com problemas graves em termos textuais e gramaticais, em relação ao que se esperaria de pessoas que estão vinculadas ao Ensino ( Técnico ) Superior.
Alguns estudantes sequer conseguiam elaborar uma frase coerente e coesa, apresentando respostas que poderiam ser classificadas como de analfabetismo funcional.
De maneira ampla, manifestou -se da parte dos estudantes uma grande dificuldade no trato da língua escrita, dificuldade essa que se traduz, sobretudo, por um desrespeito bastante frequente às regras relativas : i ) à ortografia ( por exemplo : “ proficionais ”, “ infligir as leis ”, “ corrupitos ”, “ democrácia ”, “ opitar ”, “ manter o cizilo proficionau ” ), troca de “ sociedade democrática ” por “ sociedade demográfica ” e de “ ética ” por “ étnica ” ; ii ) à pontuação ; iii ) à concordância verbal ( “ os político deve (... ) eles pode ” ) ; iv ) ao modo de organização textual ( sem distinção entre o que seja dissertativo, argumentativo ou narrativo ).
Frisa -se que esse comentário não se restringe ao que se costuma chamar de “ língua padrão ”, mas ao que poderia ser chamado de “ excelência textual ”, que não se mede exclusiva ou enfaticamente pela correção gramatical, mas pela competência, objetividade, clareza e, em alguma medida, pela criatividade no ato de redigir.
Por fim, destaca -se um dos problemas dignos de nota na formulação da questão.
A transcrição do Dicionário Houaiss apresentava defeitos.
Por exemplo, a definição 5 de “ democracia ” indicava que esse é o tipo de “ governo que acata a vontade da maioria, embora respeitando os direitos e a livre expressão das minorias ”.
O emprego de um valor concessivo na oração adverbial, por oposição a um valor aditivo, contradiz as concepções efetivamente democráticas apresentadas de 1 a 4.
Repercussões desse valor concessivo, mal empregado, foram encontradas em incontáveis respostas, as quais diziam ser típico de um governo democrático, seguir a vontade da maioria, “ apesar ” de respeitar as minorias.
Outro desses problemas dignos de nota refere -se à indicação do dicionário quanto a “ governo ”, ao passo que o comando focava a ideia mais ampla de “ sociedade democrática ”. 3.3.1.3 Análise da Questão Discursiva 10 do Componente de Formação Geral O desempenho dos estudantes na questão 10 foi inferior ao obtido na questão de número 9, com nota média 49,7 para Ingressantes, 48,8 para Concluintes.
A nota máxima também foi atingida em ambos os grupos, sendo que os percentuais de alunos que deixaram a questão em branco, tanto no grupo dos Ingressantes quanto no grupo dos Concluintes, foram ligeiramente superiores àqueles apresentados na questão 9, 16,5 \% para ambos os grupos.
Houve maior percentual de alunos, considerando -se a distribuição de notas, no intervalo entre 40 e 50 pontos, tanto para Ingressantes ( 16,7 \% ) quanto para Concluintes ( 17,0 \% ).
Também se destaca a faixa entre 70 e 80 pontos, com desempenho de Ingressantes ligeiramente superior ao dos Concluintes, respectivamente 15,9 e 15,4 \%.
Não se manifestou uma superioridade de Concluintes nas diferentes faixas de notas, com equilíbrio no desempenho dos dois grupos, ressaltando -se superioridade de Ingressantes nas faixas de notas entre 60 e 70, além da faixa seguinte, já mencionada.
Essas informações estão ilustradas na Tabela 3.14 e no Gráfico 3.11. 3.3.1.4 Comentários sobre as respostas à Questão Discursiva 10 Inicialmente, será apresentada uma panorâmica sobre as respostas, com destaque para os aspectos positivos observados.
Em relação ao desenvolvimento do tema e à vinculação do mesmo à vida cotidiana dos estudantes que participaram do exame, a visão positiva em relação à Educação de Jovens e Adultos ( EJA ) reforça a ideia de que a educação pode trazer crescimento não somente profissional, mas, sobretudo, para ampliar os horizontes da cidadania.
Foram frequentes as respostas que apontavam a educação como a melhor saída, senão a única, para o resgate da cidadania.
A atuação dos governos, em especial o governo federal, foi um ponto destacado nas respostas, as quais sinalizavam para a necessidade de as instâncias governamentais assumirem um papel cada vez mais ativo na consolidação da educação – pública, sobretudo – no país.
As políticas públicas focadas na educação foram, então, um item bastante frequente nas respostas, seja para enaltecer, criticar ou requisitar a ação governamental.
Outro ponto positivo destacado diz respeito à consciência ou ao reconhecimento da precariedade de condições de trabalho dos professores, em geral, no Brasil e, em particular, daqueles ocupados com turmas de jovens e adultos.
Nesse sentido, foram inúmeras as propostas apresentadas na segunda parte das respostas com relação à imperiosidade de melhoria salarial para os docentes, e à melhoria das condições gerais de trabalho desses profissionais.
A associação direta entre valorização do professor e sucesso escolar dos estudantes foi um ponto complementar a essa preocupação com os docentes.
Os estudantes demonstraram uma visão de vínculo no processo de ensino-aprendizagem, no qual alunos e professores são partes inalienáveis entre si.
A questão da repetência e da evasão escolar, referenciadas no suporte da questão, foram então, abordadas pelos alunos por meio do viés dessa valorização docente.
Ainda que os temas específicos da repetência e da evasão não fossem contemplados a contento na maior parte das respostas, ao menos, nessa associação com a qualificação e a valorização dos professores, algo foi trabalhado analiticamente pelos estudantes.
Num plano mais abrangente, houve reconhecimento de que, no processo ensino-aprendizagem, não há atores isolados : alunos, profissionais de educação, governos e a própria sociedade como corresponsáveis pelo sucesso ou fracasso de tal processo.
A atualidade do tema e a sua explícita vinculação com o cotidiano de muitos estudantes que realizavam o exame fizeram com que os respondentes demonstrassem um visível grau de envolvimento com o tema proposto.
Ou seja, foi possível detectar um empenho notório dos estudantes ao redigir suas respostas com seriedade e comprometimento relativo à delicada problemática dos milhões de estudantes trabalhadores do país.
Utilizar todo o espaço destinado à resposta, procurando aproveitar ao máximo a oportunidade de emitir um argumento ou uma opinião revelou -se como uma das evidências do aludido empenho.
A EJA também foi reconhecida como um dos caminhos concretos para o acesso à Universidade, no médio prazo.
As respostas se voltaram, muitas vezes, à condição de depoimentos pessoais sobre a conciliação de trabalho e estudo.
Nessa direção, ao assumirem o papel de depoentes, os estudantes expuseram suas histórias de vida, suas expectativas e percalços, suas esperanças e entraves cotidianos, seus anseios mais ambiciosos e problemas estruturais, que poderiam se resolver por meio da educação.
Nesse inciso, ficava clara a relação que se formava entre educação e mobilidade social, mediatizada por iniciativas governamentais como a EJA.
Uma sugestão plausível e válida para a solução dos problemas com a EJA no país reportou -se à necessidade de maior participação do público alvo na definição de estratégias para essa modalidade de ensino.
Nesse detalhe, era notória a associação das preocupações relativas à questão 9, que trata do problema da democracia e da cidadania, com o problema específico da educação escolar.
Houve um número significativo de sugestões ou propostas para se combater o problema da repetência e da evasão, destacando -se melhores salários para os professores e aulas mais dinâmicas.
A lista dessas sugestões, contudo, era bem mais ampla e diversificada, incluindo ainda : a redução da jornada de trabalho dos jovens e adultos que participam da EJA ; melhoria de infraestrutura nas instituições responsáveis pela EJA ; incentivo financeiro aos estudantes trabalhadores por meio de um sistema de bolsas de estudo ; realização de campanhas motivacionais para o retorno aos estudos ; estímulo à participação das empresas, por meio da solicitação de recursos materiais, espaço físico para as salas, auxílio financeiro ; incremento dos cursos profissionalizantes no país ; e ampliação da oferta de ensino a distância em todo o território nacional.
Muitos estudantes abordaram, com propriedade, a associação entre um mercado de trabalho cada vez mais competitivo e a necessidade de qualificação profissional, inclusive com uma clara percepção de que os salários são tanto maiores quanto maior a quantidade de anos de estudo e/ou a qualidade da formação.
Nesse tópico, retorna -se ao papel crucial da EJA como viabilizadora dessa formação.
Em continuidade à análise, serão apresentados aspectos negativos que foram notados com frequência nas respostas à questão 10.
De um ponto de vista mais global, a dificuldade repousa no entendimento da questão, isto é, na decifração do comando.
Muitos fugiram completamente do assunto proposto, ainda que o enunciado da questão fosse farto em dados específicos ao tema.
Em vez de responderem diretamente quais eram os problemas enfrentados por aqueles que têm de trabalhar e estudar ao mesmo tempo, muitos estudantes optavam por falar das dificuldades que existem na educação brasileira, de forma estendida e pasteurizada.
Vários estudantes não conseguiam, em suas respostas, relacionar as dificuldades de conciliação da escola com o trabalho, a não ser, como já mencionado, nos casos de depoimentos pessoais, relatos de experiência, ou seja, traços de uma biografia.
Assim, eram respostas cujo teor era apenas o de registrar o abandono dos estudos por parte de vários jovens que têm de trabalhar.
O tom das respostas, nesses casos, era meramente de uma descrição e não de uma análise, muito menos de um elemento que compusesse uma argumentação.
A dificuldade de compreensão da questão se expressava nas respostas que, em lugar de mostrarem os problemas de conciliação entre escola e trabalho, prendiam -se à ideia de “ associação ” – termo que consta do enunciado da questão – entre escola e trabalho.
Muitos alunos relatavam as consequências positivas de uma boa educação no momento em que se busca bons postos de trabalho.
Tratava -se de uma fuga parcial ao comando da questão, evidenciando má compreensão do mesmo.
Uma dificuldade facilmente percebida nas respostas refere -se à tendência de os estudantes tratarem dos problemas relativos à escola e, muito separadamente, daqueles problemas relativos ao trabalho, reconhecendo uma espécie de sinergia que impediria de visualizá -los em sua individualidade, ao mesmo tempo, que em sua complementaridade.
Isso significa que os estudantes distinguiam por mero isolamento, ou seja, eles não conseguiam distinguir e articular e, sobretudo, reconhecer as convocações mútuas entre os problemas da escola e do trabalho, no caso explícito dos estudantes trabalhadores.
Essa dificuldade de identificar, associar e analisar aspectos distintos, até mesmo concorrentes, antagônicos e complementares, denuncia a fragilidade de um pensamento simplificador que a qualidade das respostas revelou ser ainda hegemônico nas universidades brasileiras.
Notou -se a presença de preconceitos na análise dos fatores que levam ao abandono dos estudos.
Muitas respostas apontavam como causa desse abandono e da repetência a preguiça, a falta de perseverança e o cansaço dos estudantes.
Nesse tópico específico, fica evidente a responsabilização do indivíduo pelos percalços e fracassos de sua vida, descontextualizando -o de uma sociedade que, em seu movimento de conjunto, lhe impõe condicionantes.
A precariedade do uso da língua pátria de modo culto foi uma constante, conforme mencionado na análise da questão 9.
As impropriedades lexicais e a falta de elementos coesores que articulassem as partes constitutivas das respostas foram um incômodo constante.
Nesse sentido, agrava -se a situação da baixa qualidade textual em decorrência complementar da construção de textos meramente descritivos e não argumentativos, como é exigido no suporte da questão.
Isso demonstra a dificuldade dos estudantes universitários na compreensão de que uma argumentação requer, imperiosamente, justificações para os apontamentos apresentados.
Os textos das respostas eram, em sua maioria, opinativos, sem qualquer compromisso com elementos explicativos ou com a busca de justificativas. 3.3.2 Componente de Conhecimento Específico Na parte da prova relativa ao componente de Conhecimento Específico, observa -se que a diferença de desempenho foi mais acentuada entre Ingressantes e Concluintes.
A Tabela 3.15 contém dados elucidativos, na qual é possível observar que, no grupo dos Ingressantes, a nota média 33,1 é inferior à alcançada por Concluintes, 42,2.
Essas médias foram menores do que as obtidas no conjunto das questões objetivas, nas quais Ingressantes alcançaram nota média 34,1, e Concluintes, 61,8 ( Tabela 3.8 ).
O Gráfico 3.12 apresenta a distribuição de notas nos diferentes intervalos.
Observa -se que os Concluintes atingiram proporções mais elevadas na faixa entre 40 e 50 pontos, com 26,7 \% dos estudantes desse grupo situado nessa marca.
A maior concentração de Ingressantes é identificada no intervalo entre 30 e 40 pontos, correspondendo a 22,9 \% desse grupo.
A análise de cada uma destas questões será feita logo a seguir. ∗ ∗ A análise completa, com apresentação de trechos de respostas de estudantes, pode ser encontrada no documento “ Relatório da Correção das Questões Discursivas ”, elaborado pelos docentes que participaram do processo de correção das mesmas. 3.3.2.1 Análise da Questão Discursiva 38 do Componente de Conhecimento Específico Na questão 38, cujos desempenhos aferidos encontram -se na Tabela 3.16, o grupo dos Ingressantes obteve nota média 34,1, enquanto a média dos Concluintes foi 39,5.
Deixaram a questão em branco 9,2 \% dos Ingressantes e 3,9 \% dos Concluintes.
As notas dos Concluintes foram superiores às notas dos Ingressantes a partir do intervalo entre 30 e 40 pontos.
Em ambos os grupos, a nota máxima, 100,0 pontos, foi alcançada.
Uma comparação entre as três questões discursivas, com base nos dados que serão apresentados ao longo deste capítulo, demonstra que a questão 38 foi aquela em que os estudantes obtiveram o menor desempenho, dentre aquelas do Componente de Conhecimento específico. 3.3.2.2 Comentários sobre as respostas à Questão Discursiva 38 Os comentários específicos serão feitos de acordo com cada subitem que integra a questão de número 38.
Foram raras as respostas que atenderam totalmente ao previsto no padrão de resposta e que puderam ser classificadas como bem escritas.
Em relação ao item “ a ”, os estudantes mencionaram com frequência, por exemplo, a possibilidade de ovulação no período menstrual, e, nesse caso, a possibilidade da gravidez mesmo com a descamação do endométrio.
Quanto ao item “ b ”, raros foram os graduandos que mencionaram o Estatuto da Criança e do Adolescente.
Muitos deles se referiram à necessidade de o sigilo ter de ser solicitado pela adolescente, no lugar de ser um pressuposto da atividade médica, e consideraram que o médico é quem avalia se deve ou não guardar sigilo.
Em relação ao item “ c ”, muitos estudantes pareceram responder a partir de conteúdos próprios do senso comum, não mostrando domínio do conhecimento específico, ainda que com respostas corretas.
Quanto ao item “ d ”, um número expressivo de estudantes interpretou a pergunta como se fizesse referência à possibilidade de engravidar, o que se configurou como um equívoco. 3.3.2.3 Análise da Questão Discursiva 39 do Componente de Conhecimento Específico A Tabela 3.17 contém as informações relativas à questão 39 do conjunto de questões do componente de Conhecimento Específico.
O desempenho dos Concluintes, nessa questão, foi superior ao desempenho na questão de número 38, embora o mesmo não tenha acontecido com o grupo de Ingressantes.
As médias foram 32,5 no grupo dos Ingressantes e 44,9 no grupo dos Concluintes.
A proporção de estudantes que deixaram a questão em branco foi maior do que a observada na questão 38, sendo esse percentual de 21,7 \% entre os Ingressantes e de 10,8 \% dentre os Concluintes.
A nota máxima ( 100,0 ) foi atingida pelos dois grupos.
O intervalo entre 50 e 60 pontos correspondeu às notas de maior frequência para o grupo dos Concluintes e Ingressantes.
Observa -se a superioridade de Concluintes, nesta questão, em relação aos Ingressantes, especialmente após o patamar que vai de 40 a 50 pontos.
As distribuições das notas dos estudantes constam do Gráfico 3.14. 3.3.2.4 Comentários sobre as respostas à Questão Discursiva 39 Novamente, consideram os corretores, muitas respostas acompanharam o padrão do “ conhecimento geral ”, ou seja, as respostas poderiam ser formuladas por qualquer cidadão atento à publicidade ( ampla ) que é dada às medidas de controle do dengue.
A vigilância epidemiológica, em seus aspectos técnico-científicos, parece ser de desconhecimento de grande parte dos estudantes.
Os itens “ c ” e “ d ” permitiram entender melhor as deficiências de apropriação dos conteúdos específicos.
Nessa perspectiva, foram verificados alguns ‘ erros de português ’.
Além da forma, perceberam -se algumas imprecisões quanto ao significado de palavras.
Quanto ao item “ c ”, a maioria dos estudantes repetiu as principais recomendações divulgadas nos meios de comunicação.
Poucos se referiram às medidas de diagnóstico e tratamento dos pacientes.
Foi frequente a incompreensão do significado de ‘ alvenaria ’, por meio do estabelecimento de relação entre esse tipo de construção e a proliferação do mosquito.
Mesmo não tendo sido recorrente, merece destaque a \textbf{observação} de respostas com recomendação de vacinação contra dengue e de medidas de higiene pessoal como “ lavar as mãos ” e “ não levar mão suja à boca ” e também de “ uso de calçados ” e, ainda, “ identificar contactantes ”.
Também foi destacado o uso comum de “ veneno ” quando se descreve o controle das larvas.
Foi recomendado, por exemplo, o uso de ‘ agrotóxicos ’ nas residências.
Tais aspectos não foram evidenciados pelo teor pitoresco, mas sim porque revelaram que muitos dos estudantes não estão sendo atualizados quanto ao tipo de medida de controle que está sendo posto em prática pelo Ministério da Saúde.
O item “ d ” revelou o geral desconhecimento da vigilância epidemiológica, frequentemente confundida com “ vigiar a população ” no cumprimento das medidas de controle.
Diagramas de controle, níveis endêmicos, etc. foram tão raramente mencionados que é possível afirmar que tais assuntos não são do conhecimento dos estudantes.
Foi observado pelos corretores que, neste item, os estudantes deram continuidade à resposta do item anterior, acrescentando mais medidas de prevenção da proliferação do mosquito e de aparecimento de casos de dengue.
Tornou -se mais evidente, por conseguinte, o entendimento da vigilância como “ vigiar a comunidade ”, acrescido de propostas de multas e punições aos moradores que não cumprirem as determinações listadas no item “ c ”.
Como a educação médica implica capacidade de buscar educação permanente, e a prática médica envolve necessariamente a capacidade de comunicação com os pacientes e a sociedade, é preocupante a dificuldade de comunicação em língua materna demonstrada.
O processo de correção constituiu -se, por conseguinte, um indicador da necessidade de revisão de práticas de ensino para melhor formação de nossos futuros médicos. 3.3.2.5 Análise da Questão Discursiva 40 do Componente de Conhecimento Específico A Tabela 3.18 contém as informações relativas à questão 40 do conjunto do componente de Conhecimento Específico.
O desempenho dos Concluintes, nessa questão, foi inferior, ao ser comparado ao da questão 39.
Já os Ingressantes tiveram desempenho superior ao obtido na questão 38 e igual ao da questão de número 39.
A nota média de Ingressantes foi 32,5 ; já a de Concluintes, 42,1.
A nota mediana total foi 35,0.
Por outro lado, a nota máxima ( 100,0 ) foi obtida nos dois grupos.
A questão foi deixada em branco por 18,2 \% dos Ingressantes e por 4,7 \% dos Concluintes.
Os intervalos de notas nos quais foram observadas maiores frequências para Concluintes correspondem à faixa que vai de 50 a 60 pontos.
A distribuição de notas nos dois agrupamentos pode ser observada no Gráfico 3.15. observadas maiores frequências para Concluintes correspondem à faixa que vai de 50 a 60 pontos.
A distribuição de notas nos dois agrupamentos pode ser observada no Gráfico 3.15. 3.3.2.6 Comentários sobre as respostas à Questão Discursiva 40 Essa questão apresentou um número significativo de respostas em branco, possivelmente indicando a dificuldade dos estudantes em início de curso em expressar respostas diferentes das provenientes do senso comum.
Vários estudantes escreveram que ainda não tinham passado por disciplinas que permitissem responder à solicitação nela contida.
Como essa questão permitia uma resposta dissertativa, o que se constituiu em um padrão comum, as \textbf{observações} destacadas pelos corretores se referiram a pontos importantes e frequentemente observados.
Um aspecto recorrente apresentado pelos estudantes foi quanto ao tempo de \textbf{observação} do paciente.
Poucos se referiram a um tempo mínimo de \textbf{observação}, ainda que uma boa parcela tenha respondido corretamente a respeito da necessidade desse procedimento.
Um erro frequente foi a interpretação dos dados vitais da criança como “ anormais ” e não como “ normais para a idade ”.
Exames de imagem e laboratoriais ( parâmetro 1 ) foram pedidos por praticamente todos os respondentes, ainda que nem sempre de forma adequada.
A questão dos falso-negativos ( parâmetro 2 ) foi, no entanto, corretamente discutida por poucos.
Os parâmetros 3 e 4 estão relacionados, mas, como já mencionado, a indicação de necessidade de \textbf{observação} do paciente ( parâmetro 4 ) veio desacompanhada, frequentemente, da discussão do tempo necessário para descartar as possíveis complicações ( parâmetro 3 ).
O enunciado, ao induzir o estudante a selecionar a conduta de um dos médicos, produziu várias respostas do tipo “ escolho a conduta do Dr.
X ”, tornando difícil a avaliação ao se confrontar esse tipo de resposta com o gabarito.
O relato do caso e a solicitação da conduta do estudante seriam suficientes, e, ao ver dos corretores, mais adequados para a avaliação.
CONSIDERAÇÕES FINAIS A correção das questões discursivas confirmou aspectos já observados quando da correção amostral.
O processo de correção foi excelente, tendo o suporte do sistema funcionado de maneira excepcional.
Diante do volume do trabalho realizado, os contratempos foram mínimos e rapidamente sanados, deixando nos corretores impressão extremamente positiva quanto ao sistema utilizado.
Em relação à prova, a forma de apresentação das questões 38 e 39 facilitou o processo de correção.
A questão 40 permitiu uma resposta mais dissertativa, tornando a correção um pouco mais complexa.
Isso, entretanto, permitiu compreender melhor como se deu o processo de aprendizagem do estudante.

% matched lemmas: observação
\end{textsample}
